% Limitations
\section{Limitations}
\label{sec:limitations}

This work is a numerical toy demonstration with specific limitations:

\begin{itemize}
  \item \textbf{Toy landscape.} The potential is one-dimensional with six
        wells. Realistic inflationary landscapes involve many fields and
        complex tunneling dynamics.
  \item \textbf{Assigned physics.} Vacuum physics vectors
        $\boldsymbol{\theta}(v)$ are assigned, not derived from quantum
        gravity or even from the potential shape.
  \item \textbf{Simplified $\theta$ vectors.} Each basin has exactly one
        physics vector; there is no intra-basin variation.
  \item \textbf{Ancestry proxy.} The \texttt{ancestry\_proxy} kernel uses
        trajectory index as a proxy for causal ancestry. This is not true
        branch ancestry and should be treated as diagnostic only.
  \item \textbf{Diagnostic-only kernels.} \texttt{same\_basin} is
        diagnostic-only because it trivially conditions on $x_0$ basin
        membership.
  \item \textbf{Kernel and $\ell$ dependence.} Results depend on the choice
        of kernel and scale parameter $\ell$. We report sensitivity
        (Experiment~\ref{exp:4}) but do not provide a first-principles derivation
        of the optimal kernel.
  \item \textbf{Numerical, not analytic.} The observed stability improvement
        is an empirical property of this simulation, not a proven theorem
        about locality-weighted measures.
  \item \textbf{No physical claim about multiverse reality.} This simulation
        does not establish the existence or detectability of pocket universes.
  \item \textbf{Configuration-space locality as a modeling prior.} The locality
        assumption (sectors near $x_0$ in field space are more relevant) is a
        modeling prior, not an observable causal relation. It is justified only
        inside a shared-landscape model.
\end{itemize}

\ISR should be judged by stability, robustness, and failure transparency,
not by whether it can be tuned to prefer a desired vacuum.
